%% 5-11 積分は微分の逆（第11回：ページ制限なし。X16 の6・X19 に対応して書き足した）
\bkchapter{5}{11}{積分は微分の逆}{\bklevel{3}}{積分は微分の逆}
\bkmichishirube{5-4,5-9,5-10}{5-12,5-13,6-3}

\begin{goal}
\begin{itemize}[label={\color{bkTheme}●}, leftmargin=1.6\zw, itemsep=0pt, topsep=1pt]
\item \bklevel{3}\enspace 「面積がふえる速さは、その場所の高さ」であることを説明できる。
\item \bklevel{3}\enspace 原始関数を求め、不定積分を $+C$ をつけて書ける。
\item \bklevel{3}\enspace 原始関数を使って、定積分を $F(b)-F(a)$ で計算できる。
\end{itemize}
\end{goal}

\begin{hatena}
\noindent\begin{minipage}[c]{0.56\linewidth}\raggedright
5-10で、$y=x^2$ の下の面積を、細い長方形をたくさんたして求めました。0から $b$ までの面積は $\dfrac{b^3}{3}$ でした。

ところで、$(x^3)'=3x^2$（5-4）なので、$\dfrac{x^3}{3}$ を微分すると $\dfrac{3x^2}{3}=x^2$ です。面積の式を微分したら、もとのグラフの式 $x^2$ にもどりました。これは偶然でしょうか。
\end{minipage}\hfill
\begin{minipage}[c]{0.4\linewidth}
\begin{tikzpicture}[x=1mm, y=1mm, m/.style={draw=bkTheme, line width=1pt, fill=white, rounded corners=2mm, minimum width=24mm, minimum height=10mm, font=\sffamily\small}]
  \node[m] (a) at (0,12) {$x^2$};
  \node[m] (b) at (0,-12) {$\dfrac{x^3}{3}$};
  \draw[bk arrow, bkBlue] ([xshift=-4mm]a.south) -- node[left, font=\sffamily\footnotesize] {面積（積分）} ([xshift=-4mm]b.north);
  \draw[bk arrow, bkOrange] ([xshift=4mm]b.north) -- node[right, font=\sffamily\footnotesize] {微分} ([xshift=4mm]a.south);
\end{tikzpicture}
\end{minipage}
\end{hatena}

5-9の自転車でも、面積をつみ上げて位置のグラフをつくると、その傾きがもとの速さになりました。偶然ではなさそうです。この章では、「積分（面積）と微分（傾き）は逆のはたらき」であることを確かめます。これがわかると、細い長方形をたさなくても、微分の表を逆に読むだけで積分できるようになります。

\clearpage
%% ---------------- やってみよう・思い出そう ----------------
\begin{yattemiyou}[面積のふえ方を調べよう]
$y=2x$ のグラフの、0から $x$ までの面積を $S(x)$ とします。グラフがまっすぐなので、面積は三角形で求められます。
\begin{enumerate}[label={\bkcircled[bkGreen]{\arabic*}}, leftmargin=2.4\zw, itemsep=1pt, topsep=3pt]
\item $x=1$、2、3、4 のときの $S(x)$ を求め、表にしよう。
\item $x$ が1ふえるごとに、$S(x)$ がどれだけふえるかを調べよう。下の図の色の帯の面積だよ。
\item $S(x)$ を式で表して微分しよう。何になるかな？
\end{enumerate}
\begin{bkcenter}
\begin{bktable}{colspec={Q[l,m,wd=46mm] *{5}{Q[c,m,wd=9mm]}}, rowsep=1.5pt}
$x$ & 0 & 1 & 2 & 3 & 4 \\
高さ $2x$ & 0 & 2 & 4 & 6 & 8 \\
面積 $S(x)=\frac{1}{2}\times x\times2x$ & 0 & 1 & 4 & 9 & 16 \\
1ふえたときの面積のふえ方 & & 1 & 3 & 5 & 7 \\
\end{bktable}
\end{bkcenter}
\begin{bkfigure}
\begin{tikzpicture}
\begin{axis}[bkplot, width=72mm, height=50mm, xmin=0, xmax=4.6, ymin=0, ymax=9.2, xtick={1,2,3,4}, ytick={2,4,6,8}, xlabel={$x$}, ylabel={$y$}]
  \fill[bkGreen!35] (axis cs:0,0) -- (axis cs:1,0) -- (axis cs:1,2) -- cycle;
  \fill[bkBlue!30] (axis cs:1,0) -- (axis cs:2,0) -- (axis cs:2,4) -- (axis cs:1,2) -- cycle;
  \fill[bkPurple!30] (axis cs:2,0) -- (axis cs:3,0) -- (axis cs:3,6) -- (axis cs:2,4) -- cycle;
  \fill[bkOrange!35] (axis cs:3,0) -- (axis cs:4,0) -- (axis cs:4,8) -- (axis cs:3,6) -- cycle;
  \addplot[bkcurve, domain=0:4.3] {2*x};
  \node[font=\sffamily\small] at (axis cs:0.72,0.5) {1};
  \node[font=\sffamily\small] at (axis cs:1.5,1.5) {3};
  \node[font=\sffamily\small] at (axis cs:2.5,2.5) {5};
  \node[font=\sffamily\small] at (axis cs:3.5,3.5) {7};
  \node[font=\sffamily\small, text=bkBlue!80!black, anchor=west] at (axis cs:4.3,8.6) {$y=2x$};
\end{axis}
\end{tikzpicture}
\bkfigcaption{1ずつの帯。帯は台形（はじめは三角形）で、面積は1、3、5、7。}
\end{bkfigure}
\end{yattemiyou}

\noindent 三角形の底辺は $x$、高さは $2x$ なので、$S(x)=\dfrac{1}{2}\times x\times2x=x\times x=x^2$ です（$\frac{1}{2}\times2=1$）。微分すると $S'(x)=2x$（5-4）で、\textbf{グラフの高さ $2x$ と同じ}になりました。

帯で見ると、たとえば $x=1$ から2までの帯は、左の高さ2・右の高さ4の台形なので、面積は $(2+4)\div2\times1=3$ です。面積のふえ方は、その帯の高さ（の平均）と同じです。高い所ほど、面積は速くふえています。

\begin{omoidasou}
\bkrecall{5-4}{$(x^n)'=nx^{n-1}$。定数を微分すると0。定数倍・和は1つずつ微分できる。}
\bkrecall{5-10}{$\displaystyle\int_a^b f(x)\,\dd x$ は、$a$ から $b$ までの、グラフと横軸のあいだの面積（下はマイナス）。$\displaystyle\int_0^b x^2\,\dd x=\frac{b^3}{3}$}
\bkrecall{5-9}{$v$-$t$ グラフの面積をつみ上げると位置のグラフになり、その傾きはもとの速さ。}
\bkrecall{5-3}{$f'(x)=\displaystyle\lim_{h\to0}\frac{f(x+h)-f(x)}{h}$}
\bkrecall{3-10}{かけて0なら、どちらかが0。$x(x-2)=0$ なら $x=0$ か $x=2$。}
\bkrecall{2-8}{2乗して4になる数は、2と $-2$ の2つ。}
\end{omoidasou}

%% ---------------- 数学の道具 ----------------
\bksection{doguu}{数学の道具}

\bksubsection{1}{右はしを動かせる面積 $S(x)$}

5-10では、「0から1まで」「0から2まで」のように、右はしを決めて面積を求めました。こんどは、右はしを動かして、面積がどう変わるかを見ます。カーテンを右へ引いていくと、かくれる部分がだんだん広がるようすを思いうかべましょう。

\begin{bkfigure}
\begin{tikzpicture}
\pgfplotsset{mini/.style={bkplot, width=44mm, height=40mm, xmin=0, xmax=2.4, ymin=0, ymax=4.6, xtick={1,2}, ytick={2,4}, xlabel={$t$}, ylabel={$y$}}}
\begin{axis}[mini, name=m1, title={\sffamily\small 右はし $x=1$}]
  \addplot[fill=bkBlue!30, draw=none, domain=0:1] {x^2} \closedcycle;
  \addplot[bkcurve, line width=1.2pt, domain=0:2.1] {x^2};
\end{axis}
\begin{axis}[mini, name=m2, at={(m1.east)}, anchor=west, xshift=10mm, title={\sffamily\small 右はし $x=1.5$}]
  \addplot[fill=bkBlue!30, draw=none, domain=0:1.5] {x^2} \closedcycle;
  \addplot[bkcurve, line width=1.2pt, domain=0:2.1] {x^2};
\end{axis}
\begin{axis}[mini, name=m3, at={(m2.east)}, anchor=west, xshift=10mm, title={\sffamily\small 右はし $x=2$}]
  \addplot[fill=bkBlue!30, draw=none, domain=0:2] {x^2} \closedcycle;
  \addplot[bkcurve, line width=1.2pt, domain=0:2.1] {x^2};
\end{axis}
\end{tikzpicture}
\bkfigcaption{グラフ $y=t^2$ の、0から右はし $x$ までの面積。右はしを動かすと、面積が変わる。}
\end{bkfigure}

5-10の $\dfrac{b^3}{3}$ で、面積を計算できます。

\begin{bkcenter}
\begin{bktable}{colspec={Q[l,m,wd=40mm] *{5}{Q[c,m,wd=13mm]}}, rowsep=2pt}
右はし $x$ & 0 & 0.5 & 1 & 1.5 & 2 \\
面積 $S(x)=\dfrac{x^3}{3}$ & 0 & 0.0417 & 0.3333 & 1.125 & 2.6667 \\
\end{bktable}
\bktabcaption{右はし $x$ を決めると、面積が1つに決まる。だから、面積は $x$ の関数。}
\end{bkcenter}

右はし $x$ を1つ決めると、面積が1つに決まります。だから、面積は $x$ の関数です。これを $S(x)$ と書きます。

\noindent{\sffamily\bfseries\color{bkTheme}グラフの横軸を $t$ と書くわけ}\par
文字 $x$ は、もう「右はしの場所」に使いました。グラフの横軸にも $x$ を使うと、どちらの $x$ かわからなくなります。そこで、グラフの横軸は $t$ と書き、グラフを $y=t^2$ と書きます。文字がちがうだけで、5-10の $y=x^2$ と同じグラフです。定積分の記号で書くと、
\[ S(x)=\int_0^x t^2\,\dd t \]
です。「$t$ を0から $x$ まで細かく分けてたした面積」と読みます。

\bksubsection{2}{面積がふえる速さは、その場所の高さ}

\label{bk:5-11:kihon}右はしを $x$ から少しだけ動かして $x+h$ にすると、面積は少しふえます。1あたりのふえ方 $\dfrac{S(x+h)-S(x)}{h}$ は、5-2・5-3で速さを求めたときの「平均の速さ」と同じ形です。$h$ を0に近づけた先が、面積のふえる速さ $S'(x)$ です。

\noindent{\sffamily\bfseries\color{bkTheme}数で確かめる}\par
$y=t^2$、右はし $x=1$ で計算します。$S(x)=\dfrac{x^3}{3}$ を使います。たとえば $h=0.1$ なら、$S(1.1)=\dfrac{1.1^3}{3}=\dfrac{1.331}{3}$、$S(1)=\dfrac{1}{3}$ なので、ふえた面積は $\dfrac{1.331-1}{3}=0.110333\cdots$ です。これを $h=0.1$ でわると 1.10333 です。

\begin{bkcenter}
\begin{bktable}{colspec={Q[c,m,wd=12mm] Q[c,m,wd=36mm] Q[c,m,wd=24mm] Q[c,m,wd=30mm]}, rowsep=2pt, row{1}={font=\sffamily\bfseries\small}}
$h$ & ふえた面積 $S(1+h)-S(1)$ & 1あたり（$\div h$） & 右はしの高さ $(1+h)^2$ \\
0.1 & 0.110333 & 1.10333 & 1.21 \\
0.01 & 0.0101003 & 1.01003 & 1.0201 \\
0.001 & 0.0010010003 & 1.0010003 & 1.002001 \\
\end{bktable}
\bktabcaption{$h$ を小さくすると、1あたりのふえ方は1に近づく。1は、$x=1$ でのグラフの高さ $1^2$。}
\end{bkcenter}

1あたりのふえ方は、1.10333、1.01003、1.0010003と、1に近づきます。1は、$x=1$ でのグラフの高さ $1^2=1$ です。「面積がふえる速さ」が「その場所の高さ」になりました。

\noindent{\sffamily\bfseries\color{bkTheme}図で確かめる}\par
グラフ $y=f(t)$ の、$a$ から $x$ までの面積を $S(x)$ とします。右はしを $x+h$ にすると、面積は、はば $h$ の細い帯の分だけふえます。

\noindent\begin{minipage}[c]{0.5\linewidth}
\begin{tikzpicture}
\begin{axis}[bkplot, width=62mm, height=46mm, xmin=0, xmax=3.4, ymin=0, ymax=3.4, xtick={0.5,2,2.4}, xticklabels={$a$,$x$,$x\!+\!h$}, ytick=\empty, xlabel={$t$}, ylabel={$y$}]
  \addplot[fill=bkBlue!18, draw=none, domain=0.5:2] {0.6+0.35*x^2} \closedcycle;
  \addplot[fill=bkOrange!45, draw=none, domain=2:2.4] {0.6+0.35*x^2} \closedcycle;
  \addplot[bkcurve, domain=0:2.75] {0.6+0.35*x^2};
  \node[font=\sffamily\small] at (axis cs:1.3,0.5) {$S(x)$};
  \draw[bkGray, line width=0.5pt, dashed] (axis cs:2,0) -- (axis cs:2,2.0);
  \node[font=\sffamily\footnotesize, anchor=west] at (axis cs:2.5,1.0) {高さ $f(x)$};
\end{axis}
\end{tikzpicture}
\end{minipage}\hfill
\begin{minipage}[c]{0.46\linewidth}
帯（オレンジ）は、はば $h$、高さほぼ $f(x)$ の長方形なので、
\[ S(x+h)-S(x)\approx f(x)\times h \]
両辺を $h$ でわると
\[ \frac{S(x+h)-S(x)}{h}\approx f(x) \]
\end{minipage}

帯を大きくして見ましょう。帯は、2つの長方形にはさまれています（5-10の「はさみうち」と同じ考え方）。

\begin{bkfigure}
\begin{tikzpicture}[x=1mm, y=1mm]
  % 帯の拡大：左の高さ f(x)=14、右の高さ f(x+h)=24、はば h=24mm
  \fill[bkOrange!40] (0,0) -- (0,14) .. controls (9,17) and (16,20) .. (24,24) -- (24,0) -- cycle;
  \draw[bkBlue!70!black, line width=0.6pt, fill=bkBlue!18, fill opacity=0.7] (0,0) rectangle (24,14);
  \draw[bkRed!80!black, line width=0.8pt, dash pattern=on 2pt off 1.2pt] (0,0) rectangle (24,24);
  \draw[bkBlue!80!black, line width=1.4pt] (-6,10.5) .. controls (-3,12.4) .. (0,14) .. controls (9,17) and (16,20) .. (24,24) .. controls (27,25.6) .. (30,27.5);
  \draw[bkInk, line width=0.6pt] (-8,0) -- (34,0);
  \node[font=\footnotesize, anchor=north, inner sep=1pt] at (0,-1) {$x$};
  \node[font=\footnotesize, anchor=north, inner sep=1pt] at (24,-1) {$x+h$};
  \draw[bkGray, line width=0.5pt, <->] (-3,0) -- (-3,14);
  \node[font=\footnotesize, anchor=east, inner sep=1pt] at (-3.6,7) {$f(x)$};
  \draw[bkGray, line width=0.5pt, <->] (27,0) -- (27,24);
  \node[font=\footnotesize, anchor=west, inner sep=1pt] at (27.6,12) {$f(x+h)$};
  \node[font=\sffamily\footnotesize, anchor=west, align=left, inner sep=0pt] (t1) at (44,22) {{\color{bkRed!80!black}点線の長方形}：$f(x+h)\times h$（大きめ）};
  \node[font=\sffamily\footnotesize, anchor=west, align=left, inner sep=0pt] (t2) at (44,14) {{\color{bkOrange!80!black}帯}：$S(x+h)-S(x)$（ほんとう）};
  \node[font=\sffamily\footnotesize, anchor=west, align=left, inner sep=0pt] (t3) at (44,6) {{\color{bkBlue!70!black}青の長方形}：$f(x)\times h$（小さめ）};
\end{tikzpicture}
\bkfigcaption{グラフが右上がりのとき、帯は、青の長方形より大きく、点線の長方形より小さい。}
\end{bkfigure}

だから、
\[ f(x)\times h<S(x+h)-S(x)<f(x+h)\times h \]
です。全体を $h$ でわると（$h$ はプラスの数なので、大小は変わりません）、
\[ f(x)<\frac{S(x+h)-S(x)}{h}<f(x+h) \]
上の表の $h=0.1$ なら、$1<1.10333<1.21$ です。$h$ を0に近づけると、右はしの $f(x+h)$ は $f(x)$ に近づきます。まん中の数は、$f(x)$ と、$f(x)$ に近づく数にはさまれているので、$f(x)$ に近づきます。左がわは $S(x)$ の微分 $S'(x)$ の定義（5-3）そのものなので、
\[ S'(x)=f(x) \]
\textbf{面積 $S(x)$ を微分すると、グラフの高さ $f(x)$ になる}。これを、\textbf{微分積分学の基本定理}といいます。面積（積分）と傾き（微分）が、逆のはたらきでつながっていることを表しています。\bkidx{びぶんせきぶんがくのきほんていり}{微分積分学の基本定理}

\noindent{\sffamily\bfseries\color{bkTheme}たとえ}\par
へいのペンキぬりを考えます。へいの高さが場所によってちがうとき、横に\qty{1}{cm}ぬり進むごとにふえるぬった面積は、その場所のへいの高さ×\qty{1}{cm}です。高い所では面積が速くふえ、低い所ではゆっくりふえます。

\noindent{\sffamily\bfseries\color{bkTheme}物理でいうと}\par
$v$-$t$ グラフの面積は、進んだ距離でした（5-9）。0秒から $t$ 秒までの面積は、$t$ 秒での位置（出発点から）です。位置のふえる速さは、その時刻の速さです。「面積を微分すると高さ」は、「位置を微分すると速度」（5-5）と同じことを言っています。

\begin{seikaku}
上の「$f(x)$ と $f(x+h)$ にはさまれる」は、グラフが右上がりのときの話です。右下がりのときは大小が入れかわるだけで、同じように $f(x)$ に近づきます。$h$ を0に近づけると $f(x+h)$ が $f(x)$ に近づくには、$f$ がつながったグラフ（とちゅうで切れていない）であることが必要です。この巻であつかう関数は、どれもこの条件をみたしています。
\end{seikaku}

\bksubsection{3}{微分すると $f(x)$ になる関数：原始関数と不定積分}

\label{bk:5-11:genshi}②から、面積 $S(x)$ は「微分すると $f(x)$ になる関数」です。だから、面積を求めるには、微分すると $f(x)$ になる関数をさがせばよいのです。

\noindent{\sffamily\bfseries\color{bkTheme}さがしてみる}\par
微分すると $x^2$ になる関数をさがします。微分すると指数が1へるので（5-4）、指数が1大きい $x^3$ をためします。

\begin{bkcenter}
\begin{bktable}{colspec={Q[c,m,wd=30mm] Q[c,m,wd=30mm] Q[l,m,wd=52mm]}, rowsep=2pt, row{1}={font=\sffamily\bfseries\small}}
ためす関数 & 微分すると & どうする？ \\
$x^3$ & $3x^2$ & $x^2$ の3倍。多すぎるので3でわる \\
$\dfrac{x^3}{3}$ & $\dfrac{3x^2}{3}=x^2$ & ぴったり！ \\
\end{bktable}
\end{bkcenter}

このように、微分すると $f(x)$ になる関数 $F(x)$ を、$f(x)$ の\textbf{原始関数}（げんしかんすう）といいます。「原始」は「もとの」という意味で、微分する前のもとの関数ということです。$x^2$ の原始関数は $\dfrac{x^3}{3}$ です。\bkidx{げんしかんすう}{原始関数}

\noindent{\sffamily\bfseries\color{bkTheme}原始関数は1つではない}\par
$\dfrac{x^3}{3}+1$ や $\dfrac{x^3}{3}-1$ も、微分すると $x^2$ になります。定数は微分すると0だからです（5-4）。グラフで見ると、上下にずらしただけなので、同じ $x$ での傾きは変わりません。

\begin{bkfigure}
\begin{tikzpicture}
\begin{axis}[bkplot, width=84mm, height=58mm, xmin=-0.2, xmax=2.2, ymin=-1.6, ymax=3.9, xtick={1,2}, ytick={1,2,3}, xlabel={$x$}, ylabel={$y$}]
  \addplot[bkcurve, domain=-0.2:1.95] {x^3/3+1};
  \addplot[bkcurve2, domain=-0.2:1.95] {x^3/3};
  \addplot[bkcurve3, domain=-0.2:1.95] {x^3/3-1};
  \foreach \c in {1,0,-1} {\edef\tmp{\noexpand\addplot[bkInk, line width=0.9pt, domain=0.55:1.45] {1/3+\c+1*(x-1)};}\tmp}
  \addplot[only marks, mark=*, mark size=1.5pt, bkInk] coordinates {(1,4/3) (1,1/3) (1,-2/3)};
  \node[font=\footnotesize, text=bkBlue!80!black, anchor=west] at (axis cs:1.97,3.55) {$\frac{x^3}{3}+1$};
  \node[font=\footnotesize, text=bkOrange!80!black, anchor=west] at (axis cs:1.97,2.5) {$\frac{x^3}{3}$};
  \node[font=\footnotesize, text=bkGreen!60!black, anchor=west] at (axis cs:1.97,1.5) {$\frac{x^3}{3}-1$};
\end{axis}
\end{tikzpicture}
\bkfigcaption{上下にずらした3つのグラフ。$x=1$ での接線（黒）は、どれも傾き $1^2=1$ で平行。}
\end{bkfigure}

そこで、原始関数をまとめて、どれでもよい定数を $C$ として
\[ \int f(x)\,\dd x=F(x)+C\qquad（C\text{ は定数}） \]
と書き、\textbf{不定積分}（ふていせきぶん）、$C$ を\textbf{積分定数}といいます。定積分とちがって、$\int$ に下端と上端がついていません。$C$ の分だけ決まらない（定まらない）ので「不定」積分といいます。$C$ の意味は、5-13でくわしく学びます。\bkidx{ふていせきぶん}{不定積分}\bkidx{せきぶんていすう}{積分定数 $C$}

\noindent{\sffamily\bfseries\color{bkTheme}$x^n$ の不定積分}\par
同じように、$1$、$x$、$x^3$ の原始関数もさがして、ならべてみます（$1=x^0$ です。2-7）。

\begin{bkcenter}
\begin{bktable}{colspec={Q[c,m,wd=16mm] Q[c,m,wd=24mm] Q[c,m,wd=27mm] Q[c,m,wd=30mm]}, rowsep=2pt, row{1}={font=\sffamily\bfseries\small}}
$f(x)$ & 原始関数 & 微分して確かめ & 指数 \\
$1=x^0$ & $x=\dfrac{x^1}{1}$ & $(x)'=1$ & $0\to1$、1でわる \\
$x=x^1$ & $\dfrac{x^2}{2}$ & $\dfrac{2x}{2}=x$ & $1\to2$、2でわる \\
$x^2$ & $\dfrac{x^3}{3}$ & $\dfrac{3x^2}{3}=x^2$ & $2\to3$、3でわる \\
$x^3$ & $\dfrac{x^4}{4}$ & $\dfrac{4x^3}{4}=x^3$ & $3\to4$、4でわる \\
\end{bktable}
\bktabcaption{どれも「指数を1ふやして、その数でわる」になっている。}
\end{bkcenter}

きまりが見えました。どの $n$ でも、微分して確かめられます。微分のきまり $(x^n)'=nx^{n-1}$ の $n$ を $n+1$ にかえて使うと、
\[ \left(\frac{x^{n+1}}{n+1}\right)'=\frac{(n+1)x^n}{n+1}=x^n\quad\text{なので}\quad\int x^n\,\dd x=\frac{x^{n+1}}{n+1}+C\quad（n\neq-1） \]
「指数を1ふやして、その数でわる」です。微分（指数を1へらし、もとの指数をかける）のちょうど逆です。$n=-1$（$\frac{1}{x}=x^{-1}$）のときは、指数が $-1+1=0$ になり、0でわることになるので使えません（1-14）。$\frac{1}{x}$ の積分は5-14で学びます。

\noindent{\sffamily\bfseries\color{bkTheme}定数倍と和}\par
定数倍と和は、微分と同じように1つずつ積分できます。理由は2つあります。
\begin{enumerate}[label={\bkcircled[bkTheme]{\arabic*}}, leftmargin=2.4\zw, itemsep=1pt, topsep=2pt]
\item 面積で考える：高さが3倍の長方形をたせば、和も3倍です（5-10の例題6の図）。高さ $f+g$ の長方形は、$f$ と $g$ の長方形を積み重ねたものです。
\item 微分で確かめる：$(3F)'=3F'$、$(F+G)'=F'+G'$（5-4）なので、逆向きにも成り立ちます。
\end{enumerate}

\begin{bkcenter}
\begin{bktable}{colspec={Q[c,m,wd=24mm] Q[c,m,wd=46mm] Q[c,m,wd=34mm]}, rowsep=1.5pt, row{1}={font=\sffamily\bfseries\small}}
$f(x)$ & 不定積分 $\int f(x)\,\dd x$ & 確かめ（微分する） \\
$5$ & $5x+C$ & $(5x)'=5$ \\
$6x$ & $6\times\dfrac{x^2}{2}+C=3x^2+C$ & $3\times2x=6x$ \\
$3x^2+2x$ & $x^3+x^2+C$ & $3x^2+2x$ \\
\end{bktable}
\bktabcaption{不定積分は、微分して、もとにもどるかで確かめられる。}
\end{bkcenter}

\bksubsection{4}{定積分は $F(b)-F(a)$}

\label{bk:5-11:teisekibun}原始関数を使って、定積分を計算する方法を作ります。

\noindent{\sffamily\bfseries\color{bkTheme}数で考える}\par
$y=t^2$ の、1から $x$ までの面積を $S(x)$ とします。5-10の「0からでないときは、大きい面積から小さい面積をひく」を使うと、
\[ S(x)=\frac{x^3}{3}-\frac{1^3}{3}=\frac{x^3}{3}-\frac{1}{3} \]
です。$x=1$ を入れると $\dfrac{1}{3}-\dfrac{1}{3}=0$。1から1までは、はばが0なので面積も0で、合っています。$x=2$ なら $\dfrac{8}{3}-\dfrac{1}{3}=\dfrac{7}{3}$ です。

この $S(x)$ は、原始関数 $\dfrac{x^3}{3}$ に、定数 $-\dfrac{1}{3}$ をたした形です。$-\dfrac{1}{3}$ は「左はしで面積が0になるように」決まった数です。

\noindent{\sffamily\bfseries\color{bkTheme}一般にする}\par
$a$ から $x$ までの面積 $S(x)$ も、微分すると $f(x)$ になる（②）ので、$f(x)$ の原始関数の1つです。だから、どれか1つの原始関数 $F(x)$ を使って、次のように書けます。
\begin{enumerate}[label={\bkcircled[bkTheme]{\arabic*}}, leftmargin=2.4\zw, itemsep=1pt, topsep=2pt]
\item $S(x)=F(x)+C$\quad（原始関数は、$C$ だけちがう）
\item $x=a$ では、はばが0なので面積も0。だから $0=F(a)+C$、つまり $C=-F(a)$
\item ①に入れて、$S(x)=F(x)-F(a)$
\item $x=b$ とすると、$a$ から $b$ までの面積は $S(b)=F(b)-F(a)$
\end{enumerate}
\[ \int_a^b f(x)\,\dd x=F(b)-F(a)=\Bigl[\,F(x)\,\Bigr]_a^b \]
$\Bigl[\,F(x)\,\Bigr]_a^b$ は「$F(b)-F(a)$ を計算する」という書き方です。\textbf{上端を入れた値から、下端を入れた値をひきます}（「あと${}-{}$まえ」、2-3）。

\noindent{\sffamily\bfseries\color{bkTheme}物理でいうと}\par
$v$-$t$ グラフの面積は、位置の変化（変位）でした（5-9）。位置 $x(t)$ は、微分すると速度 $v$ になる関数、つまり $v$ の原始関数です。だから
\[ \int_a^b v\,\dd t=x(b)-x(a)=\text{（あとの位置）}-\text{（まえの位置）} \]
「速度の面積は、あとの位置からまえの位置をひいたもの」。上の式と同じことです。

\noindent{\sffamily\bfseries\color{bkTheme}どの原始関数を使っても、答えは同じ}\par
$x^2$ の原始関数として $\dfrac{x^3}{3}+5$ を使っても、
\[ \Bigl(\frac{8}{3}+5\Bigr)-\Bigl(\frac{1}{3}+5\Bigr)=\frac{8}{3}-\frac{1}{3}=\frac{7}{3} \]
です。$+5$ はひき算で消えます。だから、定積分では $C$ を書かなくてよいのです。

\noindent\textbf{例}\quad $\displaystyle\int_0^1x^2\,\dd x=\Bigl[\,\frac{x^3}{3}\,\Bigr]_0^1=\frac{1^3}{3}-\frac{0^3}{3}=\frac{1}{3}-0=\frac{1}{3}$。5-10で、細い長方形を1000本もたして見つけた $\dfrac{1}{3}$ が、1行で求められました。

\noindent{\sffamily\bfseries\color{bkTheme}下端と上端を入れかえると}\par
$\displaystyle\int_b^a f(x)\,\dd x=F(a)-F(b)=-\bigl(F(b)-F(a)\bigr)$ です。入れかえると符号が反対になります。5-10で「はばがマイナスになるから」と考えたことと、同じ結果です。

\bksubsection{5}{横軸より下の部分}

\label{bk:5-11:shita}定積分では、横軸より下の部分はマイナスとして計算されます（5-10）。$y=x^2-2x$ で調べます。$x^2-2x=x(x-2)$ なので、$x=0$ と $x=2$ でグラフは横軸と交わります（3-10）。

\noindent\begin{minipage}[c]{0.48\linewidth}
\begin{tikzpicture}
\begin{axis}[bkplot, width=62mm, height=62mm, xmin=-0.3, xmax=3.8, ymin=-1.6, ymax=3.6, xtick={2,3}, ytick={1,2,3}, xlabel={$x$}, ylabel={$y$}]
  \addplot[fill=bkRed!25, draw=none, domain=0:2] {x^2-2*x} \closedcycle;
  \addplot[fill=bkBlue!25, draw=none, domain=2:3] {x^2-2*x} \closedcycle;
  \addplot[bkcurve, domain=-0.25:3.25] {x^2-2*x};
  \node[font=\sffamily\small] at (axis cs:1,-0.5) {$-\frac{4}{3}$};
  \node[font=\sffamily\small, anchor=west] at (axis cs:3.2,1.1) {$+\frac{4}{3}$};
  \draw[bkGray, line width=0.4pt] (axis cs:3.22,1.0) -- (axis cs:2.85,0.55);
\end{axis}
\end{tikzpicture}
\end{minipage}\hfill
\begin{minipage}[c]{0.48\linewidth}
赤い部分は横軸より下、青い部分は横軸より上です。原始関数は $\dfrac{x^3}{3}-x^2$ なので、それぞれの定積分は次のようになります。
\end{minipage}
\begin{align*}
\int_0^2(x^2-2x)\,\dd x&=\Bigl(\frac{8}{3}-4\Bigr)-0=-\frac{4}{3}&&\text{（赤）}\\
\int_2^3(x^2-2x)\,\dd x&=(9-9)-\Bigl(-\frac{4}{3}\Bigr)=\frac{4}{3}&&\text{（青）}
\end{align*}

だから、0から3まで通して計算すると、
\[ \int_0^3(x^2-2x)\,\dd x=\Bigl[\,\frac{x^3}{3}-x^2\,\Bigr]_0^3=(9-9)-0=0 \]
となります。でも、面積が0なのではありません。下の部分（$-\frac{4}{3}$）と、上の部分（$+\frac{4}{3}$）が打ち消しあったのです。物理では、これが「もどった分を差し引いた変位」にあたります（5-9）。囲まれた部分の面積を知りたいときは、横軸と交わる所で分けて、それぞれの大きさをたします（$\frac{4}{3}+\frac{4}{3}=\frac{8}{3}$）。

\bksubsection{6}{2つのグラフをならべて見る}

\begin{bkfigure}
\begin{tikzpicture}
\begin{axis}[bkplot, name=u1, width=60mm, height=37mm, xmin=0, xmax=2.4, ymin=0, ymax=4.6, xtick={1,1.5,2}, ytick={1,2.25,4}, xlabel={$x$}, ylabel={$y$}, title={\sffamily\small 高さ\enspace$f(x)=x^2$}]
  \addplot[fill=bkBlue!18, draw=none, domain=0:1.5] {x^2} \closedcycle;
  \addplot[bkcurve, domain=0:2.1] {x^2};
  \draw[bkOrange, line width=1.2pt] (axis cs:1.5,0) -- (axis cs:1.5,2.25);
  \node[font=\sffamily\footnotesize, text=bkOrange!80!black] at (axis cs:1.8,0.45) {高さ};
  \node[font=\sffamily\footnotesize] at (axis cs:1.15,0.3) {面積};
\end{axis}
\begin{axis}[bkplot, at={(u1.east)}, anchor=west, xshift=14mm, width=60mm, height=37mm, xmin=0, xmax=2.4, ymin=0, ymax=3, xtick={1,1.5,2}, ytick={1,2}, xlabel={$x$}, ylabel={$y$}, title={\sffamily\small 面積\enspace$S(x)=\dfrac{x^3}{3}$}]
  \addplot[bkcurve2, domain=0:2.05] {x^3/3};
  \addplot[bkcurve, line width=1pt, domain=1.05:1.95] {1.125+2.25*(x-1.5)};
  \addplot[only marks, mark=*, mark size=1.6pt, bkInk] coordinates {(1.5,1.125)};
  \node[font=\sffamily\footnotesize, align=center, text=bkBlue] at (axis cs:1.75,0.5) {傾き\\2.25};
\end{axis}
\end{tikzpicture}
\bkfigcaption{左の面積（青）を右の高さにかいた。右の傾き＝左の高さ。}
\end{bkfigure}

$x=1.5$ では、右のグラフの接線の傾き $S'(1.5)=1.5^2=2.25$ が、左のグラフの高さ $f(1.5)=2.25$ と一致しています。左のグラフが高い所ほど、右のグラフは急に上がります。左のグラフが0の所（$x=0$）では、右のグラフは平ら（傾き0）です。

\begin{taisetsu}
微分積分学の基本定理\quad 面積 $S(x)$ を微分すると高さ $f(x)$\\
不定積分\quad $\displaystyle\int x^n\,\dd x=\frac{x^{n+1}}{n+1}+C$（指数を1ふやして、その数でわる）\\
定積分\quad $\displaystyle\int_a^b f(x)\,\dd x=\Bigl[\,F(x)\,\Bigr]_a^b=F(b)-F(a)$（上端の値－下端の値）
\end{taisetsu}

\begin{kotoba}
\small
\bkword{原始関数}{微分すると $f(x)$ になる関数。}
\bkword{不定積分・積分定数}{原始関数をまとめて $F(x)+C$ と書いたもの。$C$ が積分定数。}
\bkword{微分積分学の基本定理}{面積を微分すると、もとの高さになる。積分は微分の逆。}
\end{kotoba}

%% ---------------- 例題 ----------------
\bksection{reidai}{例題でたしかめよう}

\begin{reidai}[$x^n$ と定数の不定積分]
次の不定積分を求めましょう。\quad（1）$\displaystyle\int x^3\,\dd x$\quad（2）$\displaystyle\int6x^2\,\dd x$\quad（3）$\displaystyle\int5\,\dd x$\par
\kangaekata ③の「指数を1ふやして、その数でわる」を使います。数の倍（6）はそのまま残します。5は $5x^0$ と考えます。\par
\shiki（1）\nobreak 指数3を4にして、4でわる：$\dfrac{x^4}{4}+C$\par
\hspace*{2.6\zw}（2）\nobreak $6\times\dfrac{x^3}{3}+C=2x^3+C$\par
\hspace*{2.6\zw}（3）\nobreak $5\times\dfrac{x^1}{1}+C=5x+C$\par
\bkstep[bkGreen]{\R{確}{たし}かめ}微分すると、（1）$\dfrac{4x^3}{4}=x^3$、（2）$2\times3x^2=6x^2$、（3）$5$。もとにもどります。\par
\kotae（1）\nobreak $\dfrac{x^4}{4}+C$\quad（2）\nobreak $2x^3+C$\quad（3）\nobreak $5x+C$
\end{reidai}

\begin{reidai}[項がいくつかある不定積分]
$\displaystyle\int(3x^2+2x-4)\,\dd x$ を求めましょう。\par
\kangaekata 和は1つずつ積分できます（③）。項ごとに\bkreidairef{1}と同じことをして、$+C$ は最後に1つだけつけます（$C$ をいくつたしても定数）。\par
\shiki $3x^2\to3\times\dfrac{x^3}{3}=x^3$、\quad$2x\to2\times\dfrac{x^2}{2}=x^2$、\quad$-4\to-4x$\par
\hspace*{2.6\zw}合わせて $x^3+x^2-4x+C$\par
\bkstep[bkGreen]{\R{確}{たし}かめ}$(x^3+x^2-4x+C)'=3x^2+2x-4$\par
\kotae $x^3+x^2-4x+C$
\end{reidai}

\noindent\begin{minipage}[t]{0.485\linewidth}\vspace{0pt}\raggedright
\begin{jibun}
次の不定積分を求めましょう。\quad（1）$\displaystyle\int x^4\,\dd x$\quad（2）$\displaystyle\int4x\,\dd x$\quad（3）$\displaystyle\int3\,\dd x$
\bkjans{（1）\nobreak $\dfrac{x^5}{5}+C$\quad（2）\nobreak $2x^2+C$\quad（3）\nobreak $3x+C$}
\end{jibun}
\end{minipage}\hfill
\begin{minipage}[t]{0.485\linewidth}\vspace{0pt}\raggedright
\begin{jibun}
$\displaystyle\int(x^2-4x+1)\,\dd x$ を求めましょう。
\bkjans{$\dfrac{x^3}{3}-2x^2+x+C$}
\end{jibun}
\end{minipage}

\begin{reidai}[定積分]
$\displaystyle\int_1^3(2x+1)\,\dd x$ を求めましょう。\par
\kangaekata まず原始関数を1つ見つけます。$2x\to x^2$、$1\to x$ なので $F(x)=x^2+x$（$C$ は消えるので書かない）。次に、④の $F(3)-F(1)$ を計算します。\par
\shiki $F(3)=3^2+3=12$、\quad$F(1)=1^2+1=2$\par
\hspace*{2.6\zw}$\Bigl[\,x^2+x\,\Bigr]_1^3=12-2=10$\par
\bkstep[bkGreen]{\R{確}{たし}かめ}グラフ $y=2x+1$ はまっすぐなので、1から3までは、左の高さ3・右の高さ7・はば2の台形。$\dfrac{(3+7)\times2}{2}=10$。合っています。\par
\kotae 10
\end{reidai}

\begin{reidai}[速度から進んだ距離]
速度が $v=3t^2+2$（m/s）の物体が、1秒から2秒までに進んだ距離を求めましょう。\par
\kangaekata 距離は $v$-$t$ グラフの面積（5-9）。④の定積分で求めます。速度がずっとプラスなので、変位と道のりは同じです。\par
\shiki 原始関数は $t^3+2t$\par
\hspace*{2.6\zw}$\displaystyle\int_1^2(3t^2+2)\,\dd t=\Bigl[\,t^3+2t\,\Bigr]_1^2=(8+4)-(1+2)=12-3=9$\par
\kotae \qty{9}{m}
\end{reidai}

\noindent\begin{minipage}[t]{0.485\linewidth}\vspace{0pt}\raggedright
\begin{jibun}
$\displaystyle\int_0^2(3x^2+1)\,\dd x$ を求めましょう。
\bkjans{$\Bigl[\,x^3+x\,\Bigr]_0^2=10-0=10$}
\end{jibun}
\end{minipage}\hfill
\begin{minipage}[t]{0.485\linewidth}\vspace{0pt}\raggedright
\begin{jibun}
速度が $v=10t$（m/s）の物体が、0秒から3秒までに進んだ距離を求めましょう（5-2の $x=5t^2$ とくらべよう）。
\bkjans{$\Bigl[\,5t^2\,\Bigr]_0^3=\qty{45}{m}$}
\end{jibun}
\end{minipage}

\begin{column}{積分の記号の由来}
$\int$ は、ライプニッツが17世紀に使い始めた記号です。ラテン語で「和」を意味する summa の頭文字 S を縦にのばした形で、「細い長方形 $f(x)\,\dd x$ をたし合わせる」という意味がこめられています。微分の $\dd y/\dd x$ も彼の記号です。2つの記号が「逆のはたらき」をうまく表していたことが、その後の数学と物理の発展を助けました。
\end{column}

\begin{reidai}[放物線の下の面積]
$y=4-x^2$ のグラフと横軸で囲まれた部分の面積を求めましょう。\par
\kangaekata まず、グラフが横軸と交わる所（$y=0$ になる所）を求めます。$4-x^2=0$ より $x^2=4$。2乗して4になる数は $x=2$ と $x=-2$（2-8）。そのあいだでは、たとえば $x=0$ で $y=4$ なので、横軸より上です。だから、$-2$ から2までの定積分がそのまま面積です。\par
\begin{center}
\begin{tikzpicture}
\begin{axis}[bkplot, width=52mm, height=34mm, xmin=-2.7, xmax=2.9, ymin=-0.8, ymax=4.8, xtick={-2,2}, ytick={4}, xlabel={$x$}, ylabel={$y$}]
  \addplot[fill=bkBlue!25, draw=none, domain=-2:2] {4-x^2} \closedcycle;
  \addplot[bkcurve, line width=1.2pt, domain=-2.12:2.12] {4-x^2};
  \node[font=\footnotesize, text=bkBlue!80!black, anchor=west] at (axis cs:1.15,3.4) {$y=4-x^2$};
\end{axis}
\end{tikzpicture}
\end{center}
\shiki 原始関数は $4x-\dfrac{x^3}{3}$\par
\hspace*{2.6\zw}上端：$4\times2-\dfrac{2^3}{3}=8-\dfrac{8}{3}=\dfrac{16}{3}$\par
\hspace*{2.6\zw}下端：$4\times(-2)-\dfrac{(-2)^3}{3}=-8+\dfrac{8}{3}=-\dfrac{16}{3}$\par
\hspace*{2.6\zw}$\dfrac{16}{3}-\Bigl(-\dfrac{16}{3}\Bigr)=\dfrac{32}{3}$\par
\bkstep[bkGreen]{\R{確}{たし}かめ}囲まれた部分は、はば4・高さ4の長方形（面積16）の中に入っています。$\dfrac{32}{3}=10.66\cdots$ は16より小さく、ちょうど $16\times\dfrac{2}{3}$ です（放物線の山の面積は、外の長方形の $\dfrac{2}{3}$。2000年以上前にアルキメデスが見つけました）。\par
\kotae $\dfrac{32}{3}$
\end{reidai}

\begin{jibun}
$y=3x-x^2$ のグラフと横軸で囲まれた部分の面積を求めましょう。
\bkjans{$3x-x^2=x(3-x)=0$ より $x=0,\,3$。$\Bigl[\,\dfrac{3x^2}{2}-\dfrac{x^3}{3}\,\Bigr]_0^3=\dfrac{27}{2}-9=\dfrac{9}{2}$}
\end{jibun}

\begin{reidai}[面積の式を微分する]
$S(x)=\displaystyle\int_0^x(t^2+1)\,\dd t$ のとき、$S'(2)$ を求めましょう。\par
\kangaekata ②の基本定理から、$S(x)$ を微分すると、右はし $x$ でのグラフの高さ $x^2+1$ になります。積分の計算はいりません。\par
\shiki $S'(x)=x^2+1$（$t$ を $x$ にかえる）なので、$S'(2)=2^2+1=5$\par
\bkstep[bkGreen]{\R{確}{たし}かめ}$S(x)=\Bigl[\,\dfrac{t^3}{3}+t\,\Bigr]_0^x=\dfrac{x^3}{3}+x$ を微分すると $x^2+1$。同じです。\par
\kotae 5
\end{reidai}

\begin{reidai}[横軸より下の部分：変位と道のり]
速度が $v=2t-4$（m/s）の物体について、0秒から3秒までの変位（位置の変化）と、実際に動いた道のりを求めましょう。\par
\kangaekata $2t-4=0$ になるのは $t=2$。$0<t<2$ では $v<0$（もどる向き）なので、その部分の定積分はマイナスになります（⑤）。変位はそのまま計算します。道のりは、符号の変わる $t=2$ で分けて、大きさをたします（5-9）。原始関数は $t^2-4t$ です。\par
\shiki 変位\quad $\Bigl[\,t^2-4t\,\Bigr]_0^3=(9-12)-0=-3$\par
\hspace*{2.6\zw}0〜2秒\quad $\Bigl[\,t^2-4t\,\Bigr]_0^2=(4-8)-0=-4$\par
\hspace*{2.6\zw}2〜3秒\quad $\Bigl[\,t^2-4t\,\Bigr]_2^3=-3-(-4)=1$\par
\hspace*{2.6\zw}道のりは、大きさをたして $4+1=5$\par
\kotae 変位 \qty{-3}{m}、道のり \qty{5}{m}
\end{reidai}

\noindent\begin{minipage}[t]{0.485\linewidth}\vspace{0pt}\raggedright
\begin{jibun}
$S(x)=\displaystyle\int_1^x(3t^2-t)\,\dd t$ のとき、$S'(3)$ を求めましょう。
\bkjans{$3\times9-3=24$}
\end{jibun}
\end{minipage}\hfill
\begin{minipage}[t]{0.485\linewidth}\vspace{0pt}\raggedright
\begin{jibun}
速度が $v=3-t$（m/s）の物体について、0秒から4秒までの変位と道のりを求めましょう。
\bkjans{変位 $\Bigl[\,3t-\dfrac{t^2}{2}\,\Bigr]_0^4=4$。\mbox{0〜3秒}で4.5、\mbox{3〜4秒}で $-0.5$。道のり\qty{5}{m}}
\end{jibun}
\end{minipage}

\begin{reidai}[2つのグラフにはさまれた面積]
$y=x$ と $y=x^2$ のグラフで囲まれた部分の面積を求めましょう。\par
\noindent\begin{minipage}[t]{0.64\linewidth}\raggedright
\kangaekata 2つのグラフが交わる所は、$x=x^2$ より $x^2-x=0$、$x(x-1)=0$ で、$x=0,\,1$（3-10）。そのあいだでは、たとえば $x=0.5$ で $0.5>0.25$ なので、$y=x$ のほうが上です。右の図のように、はさまれた部分を細い長方形に分けると、高さは「上－下」＝ $x-x^2$ です。だから、$x-x^2$ を0から1まで積分します。
\end{minipage}\hfill
\begin{minipage}[t]{0.33\linewidth}\centering\vspace{0pt}
\begin{tikzpicture}
\begin{axis}[bkplot, width=42mm, height=40mm, xmin=0, xmax=1.25, ymin=0, ymax=1.25, xtick={1}, ytick={1}, xlabel={$x$}, ylabel={$y$}]
  \addplot[fill=bkBlue!20, draw=none, domain=0:1, samples=40] {x} \closedcycle;
  \addplot[fill=white, draw=none, domain=0:1, samples=40] {x^2} \closedcycle;
  \fill[bkOrange!60] (axis cs:0.55,0.3025) rectangle (axis cs:0.62,0.62);
  \addplot[bkcurve, line width=1.2pt, domain=0:1.15] {x};
  \addplot[bkcurve2, line width=1.2pt, domain=0:1.1] {x^2};
  \node[font=\footnotesize, text=bkBlue!80!black, anchor=south east, inner sep=1pt] at (axis cs:0.5,0.62) {$y=x$};
  \node[font=\footnotesize, text=bkOrange!80!black, anchor=west, inner sep=1pt] at (axis cs:0.82,0.42) {$y=x^2$};
\end{axis}
\end{tikzpicture}
\end{minipage}\par
\shiki $\displaystyle\int_0^1(x-x^2)\,\dd x=\Bigl[\,\frac{x^2}{2}-\frac{x^3}{3}\,\Bigr]_0^1=\frac{1}{2}-\frac{1}{3}$\par
\hspace*{2.6\zw}$\displaystyle=\frac{3}{6}-\frac{2}{6}=\frac{1}{6}$\par
\bkstep[bkGreen]{\R{確}{たし}かめ}$y=x$ の下の三角形 $\dfrac{1}{2}$ から、$y=x^2$ の下の $\dfrac{1}{3}$ をひいても $\dfrac{1}{6}$。\par
\kotae $\dfrac{1}{6}$
\end{reidai}

\begin{jibun}
$y=2x$ と $y=x^2$ のグラフで囲まれた部分の面積を求めましょう。
\bkjans{$x=0,\,2$。$\displaystyle\int_0^2(2x-x^2)\,\dd x=4-\dfrac{8}{3}=\dfrac{4}{3}$}
\end{jibun}

%% ---------------- まちがえやすいところ・物理工学 ----------------
\bksection{tsumazuki}{まちがえやすいところをチェック}

\begin{tsumazuki}{不定積分の $+C$ を書き忘れる}
\bkng{$\int2x\,\dd x=x^2$}\bkok{$\int2x\,\dd x=x^2+C$}
$x^2+1$ や $x^2-7$ も、微分すると $2x$ になる正しい答えです。不定積分には必ず $+C$ をつけます。定積分では $C$ が消えるので、いりません。
\end{tsumazuki}

\begin{tsumazuki}{指数をふやしたのに、わるのを忘れる}
\bkng{$\int x^2\,\dd x=x^3+C$}\bkok{$\int x^2\,\dd x=\dfrac{x^3}{3}+C$}
微分して確かめると、$(x^3)'=3x^2$ で、$x^2$ の3倍になってしまいます。不定積分は、いつも微分してもとにもどるかで確かめられます。
\end{tsumazuki}

\begin{tsumazuki}{$F(b)-F(a)$ の順番を逆にする}
\bkng{$\Bigl[\,x^2\,\Bigr]_1^3=1-9=-8$}\bkok{$\Bigl[\,x^2\,\Bigr]_1^3=9-1=8$}
上端（上の数）を入れた値から、下端（下の数）を入れた値をひきます（「あと－まえ」、2-3）。
\end{tsumazuki}

\begin{tsumazuki}{下端の値が0だと思いこむ}
\bkng{$\Bigl[\,x^2+x\,\Bigr]_1^3=9+3=12$}\bkok{$\Bigl[\,x^2+x\,\Bigr]_1^3=12-2=10$}
下端が0でないときは、下端を入れた値も計算して、ひきます。下端が0でも、$F(0)$ が0とは限りません（$F(x)=x+5$ なら $F(0)=5$）。
\end{tsumazuki}

\begin{butsuri}
\begin{itemize}[leftmargin=1.2\zw, itemsep=2pt, topsep=0pt, label={\color{bkPurple}●}]
\item \textbf{ニュートンとライプニッツの発見}\par
毎回細かくたす大変な計算だった積分が、微分の逆として求められるとわかり、物理が一気に進みました。
\item \textbf{位置・速度・加速度}\par
速度を積分すると位置、加速度を積分すると速度。微分と積分を行き来して運動を調べます（5-13）。
\item \textbf{積算計}\par
ガスや水道のメーターは、流れる速さ（流量）を時間で積分して、使った量を表示しています（5-12）。
\end{itemize}
\end{butsuri}

%% ---------------- 章末問題 ----------------
\bkendproblems
\bkproblemgroup{基本}{3}

\bkproblem $\displaystyle\int(4x^3-6x+2)\,\dd x$ を求めましょう。
\bkpans{$x^4-3x^2+2x+C$}

\bkproblem $\displaystyle\int_1^2(3x^2-2x)\,\dd x$ を求めましょう。
\bkpans{$\Bigl[\,x^3-x^2\,\Bigr]_1^2=4-0=4$}

\bkproblem 速度が $v=6t-2$（m/s）の物体が、0秒から3秒までに進んだ距離を求めましょう。
\bkpans{$\Bigl[\,3t^2-2t\,\Bigr]_0^3=\qty{21}{m}$}

\bkproblem $y=4x-x^2$ のグラフと横軸で囲まれた部分の面積を求めましょう。
\bkpans{$\Bigl[\,2x^2-\dfrac{x^3}{3}\,\Bigr]_0^4=\dfrac{32}{3}$}

\bkproblem 速度が $v=2t-6$（m/s）の物体について、0秒から4秒までの変位と道のりを求めましょう。
\bkpans{変位 $\Bigl[\,t^2-6t\,\Bigr]_0^4=\qty{-8}{m}$。\mbox{0〜3秒}で $-9$、\mbox{3〜4秒}で1。道のり\qty{10}{m}}

\bkproblem $y=3x$ と $y=x^2$ のグラフで囲まれた部分の面積を求めましょう。
\bkpans{$\displaystyle\int_0^3(3x-x^2)\,\dd x=\dfrac{27}{2}-9=\dfrac{9}{2}$}

\bkproblem $\displaystyle\int(x^3+6x^2-5)\,\dd x$ を求め、微分して、もとにもどることを確かめましょう。
\bkpans{$\dfrac{x^4}{4}+2x^3-5x+C$（微分すると $x^3+6x^2-5$）}

\bkproblemgroup{チャレンジ}{3}

\bkproblem $S(x)=\displaystyle\int_0^x(t^3-2t)\,\dd t$ のとき、$S'(2)$ を求めましょう。また、$S(x)$ を計算して、微分した結果と同じになることを確かめましょう。
\bkpans{$S'(2)=8-4=4$。$S(x)=\dfrac{x^4}{4}-x^2$ を微分すると $x^3-2x$}

\bkproblem $\displaystyle\int_0^3(x^2-2x)\,\dd x=0$ ですが、グラフと横軸で囲まれた部分の面積は0ではありません。$0\leqq x\leqq2$ と $2\leqq x\leqq3$ に分けて、それぞれの定積分を求めましょう。
\bkpans{$\Bigl[\,\dfrac{x^3}{3}-x^2\,\Bigr]_0^2=-\dfrac{4}{3}$、$\Bigl[\,\dfrac{x^3}{3}-x^2\,\Bigr]_2^3=0-\left(-\dfrac{4}{3}\right)=\dfrac{4}{3}$}

\bkproblem $S(x)=\displaystyle\int_1^x(2t+3)\,\dd t$ を、$x$ の式で表しましょう。また、$S(1)=0$ と、$S'(x)=2x+3$ になることを確かめましょう。
\bkpans{$S(x)=\Bigl[\,t^2+3t\,\Bigr]_1^x=x^2+3x-4$。$S(1)=1+3-4=0$、$S'(x)=2x+3$}

\bkproblem $y=x^2-4$ のグラフと横軸で囲まれた部分の面積を求めましょう。
\bkpans{$\displaystyle\int_{-2}^{2}(x^2-4)\,\dd x=-\dfrac{32}{3}$（横軸より下なのでマイナス）。面積は $\dfrac{32}{3}$}

\begin{furikaeri}
\begin{checklist}
\item 面積のふえる速さが高さであることを説明できる。\hfill{\footnotesize\sffamily\color{bkTheme}→\pageref{bk:5-11:kihon}ページ}
\item 不定積分を $+C$ をつけて書ける。\hfill{\footnotesize\sffamily\color{bkTheme}→\pageref{bk:5-11:genshi}ページ}
\item 定積分を $F(b)-F(a)$ で計算できる。\hfill{\footnotesize\sffamily\color{bkTheme}→\pageref{bk:5-11:teisekibun}ページ}
\end{checklist}
\end{furikaeri}

\bkfukidashi{sora}{積分がかんたんになった！　次は、物理で積分を使ってみよう。}
